\chapter{Netzgrößen, Leistung und Kurzschlussleistung}
\label{chap:netzgroessen}

\begin{lernziel}
Nach diesem Kapitel kannst du:
\begin{itemize}
  \item Netz-Nennspannung, Leiter-Leiter-Spannung und Strangspannung sicher unterscheiden,
  \item Scheinleistung und Nennstrom eines Transformators beziehungsweise Netzes berechnen,
  \item den Zusammenhang zwischen Kurzschlussleistung, Netzimpedanz und prospektivem Kurzschlussstrom herleiten,
  \item die Grenzen einer vereinfachten Kurzschlussrechnung benennen,
  \item Ergebnisse auf Plausibilität und Einheiten prüfen.
\end{itemize}
\end{lernziel}

\section{Nennwerte und tatsächliche Betriebsgrößen}

Die Nennspannung \(U_\mathrm{n}\) ist eine festgelegte Kenngröße zur Bezeichnung
eines Netzes oder Betriebsmittels. Sie ist nicht automatisch identisch mit der
Spannung, die zu jedem Zeitpunkt an jedem Knoten gemessen wird. Die tatsächliche
Betriebsspannung hängt unter anderem von Lastfluss, Netzregelung, Leitungsimpedanz
und Betriebszustand ab.

Bei Drehstromnetzen bezeichnet die übliche Netzspannungsangabe die Spannung
zwischen zwei Außenleitern. Für ein symmetrisches Netz gilt:
\[
U_\mathrm{LL}=\sqrt{3}\,U_\mathrm{Str}.
\]
Bei einem Netz mit \(U_\mathrm{LL}=\SI{20}{\kilo\volt}\) beträgt die
Strangspannung im symmetrischen Zustand daher
\[
U_\mathrm{Str}=\frac{\SI{20}{\kilo\volt}}{\sqrt{3}}
\approx \SI{11.55}{\kilo\volt}.
\]
Diese Beziehung darf nicht ohne Weiteres auf unsymmetrische Fehlerzustände
mit Sternpunktverschiebung übertragen werden.

\section{Drehstromleistung und Bemessungsstrom}

Für ein symmetrisches Drehstromsystem ist die Scheinleistung
\begin{equation}
  S=\sqrt{3}\,U_\mathrm{LL} I_\mathrm{L}.
  \label{eq:netz_scheinleistung}
\end{equation}
Daraus folgt für den Leiterstrom:
\begin{equation}
  I_\mathrm{L}=\frac{S}{\sqrt{3}\,U_\mathrm{LL}}.
  \label{eq:netz_nennstrom}
\end{equation}
Die Gleichung ist beispielsweise für die Plausibilisierung von
Transformator-Bemessungsströmen hilfreich. Sie sagt allein jedoch noch nichts
über zulässige Dauerbelastung, Kühlung, Umgebungstemperatur oder Schutzkonzept
aus.

\begin{rechenbeispiel}
Ein Transformator hat eine Bemessungsscheinleistung von
\(\SI{40}{\mega\volt\ampere}\) und eine Oberspannung von
\(\SI{110}{\kilo\volt}\). Der Bemessungsstrom auf der Oberspannungsseite ist
\[
I_\mathrm{n,OS}=
\frac{\SI{40}{\mega\volt\ampere}}
{\sqrt{3}\cdot\SI{110}{\kilo\volt}}
\approx \SI{210}{\ampere}.
\]
Auf einer Unterspannungsseite mit \(\SI{20}{\kilo\volt}\) ergibt sich bei
derselben Bemessungsscheinleistung:
\[
I_\mathrm{n,US}=
\frac{\SI{40}{\mega\volt\ampere}}
{\sqrt{3}\cdot\SI{20}{\kilo\volt}}
\approx \SI{1.155}{\kilo\ampere}.
\]
Die niedrigere Spannung führt bei gleicher Leistung zu einem höheren Strom.
Das ist einer der Gründe, weshalb Übertragungsnetze hohe Spannungen verwenden.
\end{rechenbeispiel}

\section{Netzimpedanz und Kurzschlussstrom}

Ein Kurzschluss führt zu einem hohen Strom, weil die Impedanz zwischen
speisender Quelle und Fehlerort stark kleiner ist als die Impedanz im
Normalbetrieb mit angeschlossener Last. Für eine erste einphasige Näherung
wird das Netz durch eine ideale Spannungsquelle hinter einer Ersatzimpedanz
\(Z_k\) modelliert.

Für einen Fehler an einer Effektivwertquelle gilt im einfachen Thevenin-Modell:
\begin{equation}
  \phasor I_k=\frac{\phasor U_\mathrm{Th}}{Z_k}.
  \label{eq:thevenin_ik}
\end{equation}
Der Betrag des Stroms ist damit
\[
I_k=\frac{U_\mathrm{Th}}{|Z_k|}.
\]
Wichtig: Spannung und Impedanz müssen auf dieselbe Ersatzschaltung und
Spannungsebene bezogen sein. In einer dreiphasigen Kurzschlussrechnung ist
bei Verwendung der Leiter-Leiter-Spannung die Strangspannung
\(U_\mathrm{LL}/\sqrt{3}\) die treibende Spannung des einphasigen
Ersatzmodells.

\section{Kurzschlussleistung}

Die Kurzschlussleistung ist eine Kenngröße für die elektrische Stärke eines
Netzknotens. In einer üblichen dreiphasigen Definition wird sie aus der
Nennspannung und dem anfänglichen symmetrischen Kurzschlusswechselstrom
gebildet:
\begin{equation}
  S_k''=\sqrt{3}\,U_\mathrm{n}\,I_k''.
  \label{eq:sk_definition}
\end{equation}
Daraus folgt:
\begin{equation}
  I_k''=\frac{S_k''}{\sqrt{3}\,U_\mathrm{n}}.
  \label{eq:ik_from_sk}
\end{equation}
Die doppelte Strichnotation \(I_k''\) wird in der Kurzschlussberechnung
häufig für den Anfangs-Kurzschlusswechselstrom verwendet. In konkreten
Normen und Rechenverfahren ist die jeweils definierte Stromgröße zu beachten.

\begin{merke}
Die Kurzschlussleistung \(S_k''\) ist keine während des Fehlers tatsächlich
in Wärme umgesetzte Wirkleistung. Sie ist eine aus Spannung und
Kurzschlussstrom gebildete Kenngröße des Netzknotens.
\end{merke}

\begin{rechenbeispiel}
An einem \(\SI{20}{\kilo\volt}\)-Netzknoten sei die dreiphasige
Kurzschlussleistung \(S_k''=\SI{250}{\mega\volt\ampere}\) angegeben.
Der anfängliche symmetrische Kurzschlussstrom beträgt näherungsweise
\[
I_k''=\frac{\SI{250}{\mega\volt\ampere}}
{\sqrt{3}\cdot\SI{20}{\kilo\volt}}
\approx \SI{7.22}{\kilo\ampere}.
\]
Die Größenordnung ist plausibel: Eine Kurzschlussleistung von
\(\SI{250}{\mega\volt\ampere}\) bei nur \(\SI{20}{\kilo\volt}\) führt zu
einem Strom im Bereich mehrerer Kiloampere.
\end{rechenbeispiel}

\section{Abschätzung der Netzimpedanz aus der Kurzschlussleistung}

Für eine erste Näherung eines symmetrischen dreipoligen Kurzschlusses kann
der Betrag der äquivalenten Netzimpedanz auf der betrachteten Spannungsebene
abgeschätzt werden durch
\begin{equation}
  |Z_k|\approx\frac{U_\mathrm{n}^2}{S_k''}.
  \label{eq:zk_from_sk}
\end{equation}
Dabei ist \(U_\mathrm{n}\) die Leiter-Leiter-Spannung. Die Formel ergibt sich
aus dem einphasigen Ersatzmodell mit Strangspannung \(U_\mathrm{n}/\sqrt{3}\)
und dem Zusammenhang \(S_k''=\sqrt{3}U_\mathrm{n}I_k''\):
\[
I_k''\approx\frac{U_\mathrm{n}/\sqrt{3}}{|Z_k|}
\quad\Longrightarrow\quad
S_k''\approx\frac{U_\mathrm{n}^2}{|Z_k|}.
\]
Die Beziehung liefert zunächst nur den Betrag. Für die Aufteilung in
\(R_k\) und \(X_k\) ist zusätzliche Information über das Verhältnis \(R/X\)
oder eine detailliertere Netzberechnung nötig.

\begin{rechenbeispiel}
Für den zuvor betrachteten Netzknoten gilt:
\[
|Z_k|\approx
\frac{(\SI{20}{\kilo\volt})^2}{\SI{250}{\mega\volt\ampere}}
=\SI{1.6}{\ohm}.
\]
Die Einheit lässt sich direkt prüfen:
\(\si{\volt\squared}/\si{\volt\ampere}=\si{\ohm}\).
Der Wert ist eine auf die \(\SI{20}{\kilo\volt}\)-Ebene bezogene
Ersatzimpedanz für die vereinfachte dreipolige Kurzschlussbetrachtung.
\end{rechenbeispiel}

\section{Was die einfache Rechnung nicht abbildet}

Die Formeln oben sind nützlich für erste Abschätzungen, ersetzen aber keine
vollständige Kurzschlussberechnung. Je nach Aufgabe sind unter anderem zu
berücksichtigen:
\begin{itemize}
  \item Impedanzen von Transformatoren, Leitungen und weiteren Quellen,
  \item Beiträge von Motoren, Generatoren und leistungselektronischen Anlagen,
  \item das Verhältnis \(R/X\) und der Gleichstromanteil des Kurzschlussstroms,
  \item der zeitliche Verlauf und die Unterscheidung zwischen Anfangs-,
        Ausschalt- und stationärem Kurzschlussstrom,
  \item Fehlerart, Fehlerimpedanz, Erdungsart und Nullsystemimpedanz,
  \item der verwendete Berechnungsstandard und dessen Spannungsfaktoren.
\end{itemize}

Ein besonders wichtiger Punkt: Die Kurzschlussleistung allein beschreibt
nicht alle Fehlerarten. Für einen Erdschluss oder einen zweipoligen Fehler
müssen die jeweiligen Sequenzimpedanzen und Netzerdungsverhältnisse
berücksichtigt werden.

\section{Plausibilitätsprüfungen}

Vor jeder Weiterverwendung sollte man mindestens prüfen:
\begin{enumerate}
  \item Stimmen die Spannungsebene und die Bezugsseite?
  \item Wurde Leiter-Leiter-Spannung oder Strangspannung eingesetzt?
  \item Ist der berechnete Strom in der erwarteten Größenordnung?
  \item Ist \(S_k''=\sqrt{3}U_nI_k''\) mit den eingesetzten Werten erfüllt?
  \item Ist klar, ob ein Anfangs-, Ausschalt- oder Dauerstrom gemeint ist?
  \item Handelt es sich um einen dreipoligen Kurzschluss oder einen Erdschluss?
\end{enumerate}

\section{Übungsaufgaben}

\begin{aufgaben}
  \item Ein \(\SI{63}{\mega\volt\ampere}\)-Transformator hat auf der
        \(\SI{110}{\kilo\volt}\)-Seite seine Bemessungsspannung.
        Berechne den Bemessungsstrom auf dieser Seite.
  \item Ein \(\SI{10}{\kilo\volt}\)-Netzknoten hat eine angegebene
        dreiphasige Kurzschlussleistung von \(\SI{180}{\mega\volt\ampere}\).
        Berechne den anfänglichen symmetrischen Kurzschlussstrom.
  \item Bestimme aus Aufgabe 2 den Betrag der vereinfachten
        Thevenin-Ersatzimpedanz auf der \(\SI{10}{\kilo\volt}\)-Ebene.
  \item Ein Techniker setzt in einem dreiphasigen Kurzschlussmodell
        \(U_\mathrm{n}\) direkt als Strangspannung ein. Erkläre den Fehler
        und die Auswirkung auf das Ergebnis.
  \item Warum reicht die Kurzschlussleistung eines Netzknotens allein
        nicht aus, um einen einpoligen Erdschlussstrom zu bestimmen?
\end{aufgaben}

\section{Kurzlösungen}

\begin{enumerate}[label=\arabic*.]
  \item \(I_n=63/(\sqrt{3}\cdot110)\,\si{\mega\volt\ampere\per\kilo\volt}
        \approx\SI{331}{\ampere}\).
  \item \(I_k''=180/(\sqrt{3}\cdot10)\,\si{\mega\volt\ampere\per\kilo\volt}
        \approx\SI{10.39}{\kilo\ampere}\).
  \item \(|Z_k|\approx U_n^2/S_k''=(10\,\si{\kilo\volt})^2/
        180\,\si{\mega\volt\ampere}\approx\SI{0.556}{\ohm}\).
  \item Die treibende Spannung des einphasigen Ersatzmodells ist
        \(U_n/\sqrt{3}\). Eine Verwechslung um den Faktor \(\sqrt{3}\)
        führt entsprechend zu einem falschen Strom.
  \item Dafür werden insbesondere die Nullsystemimpedanz und die
        Sternpunktbehandlung benötigt; diese werden in späteren Kapiteln
        eingeführt.
\end{enumerate}